\begin{theorem}[Wild odd-prime non-stable case]\label{mod:cases-wildodd-main}
Let \(K/F\) be a wildly ramified prime-cyclic extension of
nonarchimedean local fields, with positive lower break \(t\), residue degree
one, and
\[
 [K:F]=p=\operatorname{char}(k_F),\qquad p\ne2.
\]
Fix an integral uniformizer \(\pi_K\) which generates
\(\mathcal O_K\) over \(\mathcal O_F\).  Let
\(\Delta_F,\Delta_K\) be local-constant functions satisfying
\texttt{IsDeltaFiniteLocalConstant}.

For exact multiplicative and additive data \(\chi_F,\psi_F\), let
\(\chi_{\min}\) be the least-conductor member of the complete
norm-character orbit of \(\chi_F\).  Higher data are required only under the
implication \(1<m_F(\chi_{\min})\).  In that branch,
\texttt{WildOddNonstableHigherData} supplies:
\begin{itemize}
\item the actual decomposition
  \(m_F(\chi_{\min})=2d+\varepsilon\), with its exact parity proof and the
  non-stable inequality \(d<t+1\);
\item complete \texttt{FirstMainComputationalData},
  \texttt{FirstMainPhaseData}, and the \texttt{ExactOddPhaseAssembly}
  indexed by all nonidentity norm characters at the actual conductor;
\item the certified residual additive character and the completed
  \texttt{WildOddParityReductionData}, together with the identification of
  the assembly's residual phase with the complete residual phase quotient;
\item the completed \texttt{WildOddCorrectionData} at \(T=t+1\), tied to
  the actual conductor, the assembly element \(u\), and the exact correction
  coordinate \(X\); and
\item a nontrivial norm character of conductor \(T\), with its proved
  stationary linearization on a positive depth \(r\le T\).
\end{itemize}
No stationary datum is requested when \(m_F(\chi_{\min})\le1\).

Then Lean proves the literal First Main identity
\[
 \Delta_K(\chi_F\circ N_{K/F},\psi_F\circ\operatorname{Tr}_{K/F})
 \prod_{\mu\in S(K/F)}\Delta_F(\mu,\psi_F)
 =
 \prod_{\mu\in S(K/F)}\Delta_F(\mu\chi_F,\psi_F),
\]
where both products range over the complete norm-character group, including
the identity.

The proof first reduces to \(\chi_{\min}\).  Actual conductors zero and one
invoke \texttt{firstMain\_wildEndpoint\_zero} and
\texttt{firstMain\_wildEndpoint\_one}; the finite-Delta hypotheses identify
their concrete finite sums with the displayed local constants.  In the
higher branch, \texttt{PhaseReductionResult} gives the exact equality
\[
 \operatorname{Err}_{K/F}(\chi_{\min})
 =\mathcal R_{\mathrm{odd}}\,\psi_F(-AX).
\]
The complete theorem exported by \texttt{wildOdd\_parityReduction} gives
\(\mathcal R_{\mathrm{odd}}=1\).  The correction-coordinate equality and
\texttt{wildOdd\_correction\_phase} give \(\psi_F(-AX)=1\), using the
conductor-\(T\) norm character and its supplied linearization.  Hence the
complete error term is one.

Finally, the proved equivalence between error one and
\texttt{FirstMainIdentity} is applied.  The minimal representative's exact
computational data are reindexed through its norm-character twist, and the
proved equality
\(\operatorname{Err}_{K/F}(\chi_{\min})=
  \operatorname{Err}_{K/F}(\chi_F)\)
transports the result back in the established direction.
\LeanFile{LanglandsFirstMainLemma/Cases/WildOdd/Main.lean}
\lean{LanglandsFirstMainLemma.WildOddNonstableHigherData}
\lean{LanglandsFirstMainLemma.firstMain_wildOdd_nonstable}
\leanok
\SourceText{Corollary cor:wild-odd.}
\uses{mod:cases-wildendpoints,mod:cases-wildodd-parityreduction,mod:cases-wildodd-correctionvaluation,mod:parameters-phasereduction,mod:parameters-minimalorbitstationary}
\FormalizationContract
The higher package consumes only the completed public interfaces of the five
dependency nodes.  It does not reconstruct stationary coefficients, choose
canonical representatives, transport residual coordinates, prove a new
finite-field phase identity, or establish a new norm--trace valuation.  Its
correction conductor, element, and coordinate are explicitly identified
with the actual phase assembly.

The theorem is an assembly statement for exact conductor packages and
computational local constants.  The non-stable boundary is the literal
\(d<t+1\), complementary to the stable hypothesis \(t+1\le d\).  Endpoint
conductors are dispatched before any higher package is opened.  The public
conclusion is \texttt{FirstMainIdentity}, not merely a residual phase or an
error-term equality.
\end{theorem}
